\section{Промышленный цикл как специфическая форма диалектического развития капиталистической экономики}

\subsection{Материалистическое основание цикла}
Исследование цикла должно исходить не из внешней периодичности явлений, а из развития как процесса возникновения, развертывания и снятия противоречий, определяющих последовательную смену его фаз, поскольку категория развития выступает исходной по отношению к категории цикла и позволяет раскрыть внутреннюю логику циклического движения, обусловленного имманентными противоречиями системы.

\textit{Материя} как философская категория обозначает объективную реальность, существующую независимо от сознания человека и первичную по отношению к нему, тогда как \textit{движение} выступает всеобщим способом существования материи, выражающим неотделимые от нее изменения, превращения, взаимодействия и внутренние противоречия; а \textit{развитие} свою очередь представляет собой высшую форму данного движения, характеризующуюся направленным, необратимым и закономерным изменением материальных и идеальных систем, в ходе которого совершается переход от одного качественного состояния к другому, более сложному, развитому и внутренне преобразованному, вследствие чего развитие следует понимать не как простое изменение вообще, а как движение, обладающее собственной внутренней логикой, направленностью и результатом в виде возникновения новых форм бытия.

Таким образом, \textit{цикл} выражает собой развития в его относительно завершенной, повторяющейся и вместе с тем изменяющейся форме, характеризующейся возвратом к исходным моментам не в качестве простого повтора прежнего состояния, а в качестве его преобразования на новой ступени, сочетающей в себе воспроизводство вместе с качественным изменением и переходом к новому состоянию системы.

\subsection{Системный характер цикла}
Цикл может быть понят исключительно в качестве формы движения системы, поскольку его фазы, ритм и направление определяются не случайными внешними воздействиями, а имманентной структурой элементов, отношений и противоречий системы, задающих объективные ограничения и траекторию ее развития.

Каждая \textit{система} содержит в себе заимосвязанные, взаимнообусловленные, взаимноисключающие друг друга и находящиеся в непрерывном состоянии единства и борьбы диалектические \textit{противоположности}, обуславливающие ее имманентные \textit{противоречия}, среди которых можно выделить \textit{главные}, определяющие тип, структуру, динамику и направление ее развития в целом, и \textit{второстепенные}, производные от главных и подчиненные им; в то время как цикл представляет собой закономерную форму развития системы, подразумевающую непрерывный процесс возникновения, развертывания и разрешения ее противоречий, подчиненный логике всеобщих \textit{законов материалистической диалектики}: закону единства и борьбы противоположностей, закону перехода количественных изменений в качественные и закону отрицания отрицания.

Таким образом, системный характер цикла заключается не во внешней периодичности повторов состояний, а в форме диалектического развития, подразумевающей закономерную смену фаз развития данной системы посредством разрешения присущих ей диалектических противоречий и обуславливающей тем самым воспроизводство и преобразование самой ее целостности

\subsection{Механизм диалектического развития цикла}
Применительно к циклу механизм диалектического развития выражает внутреннюю логику последовательной смены его фаз, представляя собой не простую повторяемость состояний, а последовательное накопление, обострение, разрешение и снятие противоречий системы, образующее относительно завершенный виток в рамках спиралевидной траектории ее движения и переход от одного состояния к другому.

Механизм диалектического развития системы представляет собой непрерывный процесс единства и борьбы ее противоположностей, обуславливающий переход ее количественных изменений в качественные и следующее за ним снятие исходных противоречий, находящий свое выражение в четырех последовательны этапах:
\begin{enumerate}
    \item \textit{Накопление} противоречий системы в скрытом виде посредством постепенных изменений количественного соотношения ее диалектических противоположностей;
    \item \textit{Обострение} накопленных противоречий системы в явном виде до критического уровня вследствие существенного изменения количественного соотношения ее диалектических противоположностей;
    \item \textit{Разрешение} не совместимых со существование текущей системы противоречий насильственным путем посредством разрыва диалектического единства ее противоположностей в ходе кризиса;
    \item \textit{Снятие} противоречий прежней системы в ходе ее отрицания или перехода на качественно новый уровень развития посредством формирования нового диалектического единства противоположностей.
\end{enumerate}

Таким образом, \textit{кризис} представляет собой закономерный переломный момент, выступающий конституирующей фазой циклического развития системы, выражающийся в разрыве прежней формы единства ее противоположностей и обуславливающий тем самым ее переход к качественно новому состоянию в рамках каждого последующего витка спирали диалектического развития.

\subsection{Цикл как специфическая форма диалектического развития}
Необходимость рассмотрения цикла именно как специфической формы диалектического развития проистекает не из механического характера повторения одних и тех же состояний системы, а из ее закономерного движения через последовательную смену фаз, определяемую возникновением, обострением, разрешением и последующим воспроизводством ее внутренних противоречий, способных быть должным образом раскрытыми исключительно в рамках категорий материалистической диалектики.

В общем смысле цикл представляет собой форму существования и движения системы, процесса или совокупности элементов, подразумевающую последовательный, внутренне связанный, относительно завершенный и воспроизводимый характер ее изменений, выражающийся в целостном витке развития или ряде фаз, звеньев и частей, закономерно объединенных в рамках некоторого структурного единства; отвечающую критериям:
\begin{itemize}
    \item \textit{Воспроизводимости морфологии фаз} в качестве наличия устойчивой к малым возмущениям типовой, закономерной, последовательной и распознаваемой структуры этапов развития системы;
    \item \textit{Инвариантности и преемственности} в качестве наличия обуславливающих возможность сравнительного анализа в различные периоды развития системы ее неизменных или сопоставимых признаков и параметров;
    \item \textit{Временной когерентности и регулярной периодичности} в качестве наличия обуславливающей возможность выделения характерных, сопоставимых, закономерных периодов развития системы когерентной траектории изменения ее параметров с допустимыми, не нарушающими морфологию фаз временными разрывами;
    \item \textit{Индексируемости и формализуемости} в качестве наличия обуславливающих возможность установления и обоснования в ходе практики существования или отсутствия циклических закономерностей и зависимостей развития системы количественных и качественных параметров и величин.
\end{itemize}

Таким образом, в контексте материалистической диалектики цикл можно определить в качестве специфической формы диалектического развития системы, характеризующейся временным разрешением, частичным снятием и дальнейшим воспроизводством ее преимущественно второстепенных противоречий на качественно новом уровне ее развития при отсутствии ее диалектического отрицания и сохранении ее фундаментальных противоречий в совокупности с обуславливающими их системообразующими диалектическими противоположностями.

\subsection{Капиталистическая экономика как сфера промышленного цикла}
Для целей раскрытия природы промышленного цикла необходимым является определение той исторически конкретной экономической среды, в рамках которой он возникает, развивается и приобретает закономерный характер, поскольку цикл не существует вообще, вне определенного способа производства, а всегда выражает движение присущих ему производительных сил, производственных отношений и внутренних противоречий, определяющих особенности воспроизводственного процесса, формы его неравномерности и механизмы перехода от подъема к кризису, вследствие чего анализ промышленного цикла предполагает рассмотрение капиталистической экономики не просто как совокупности хозяйственных явлений, а как особой системы общественного воспроизводства, в пределах которой циклическая динамика получает свое объективное основание.

\textit{Капиталистическая экономика} отражает исторически конкретную общественно-экономическую формацию и систему социальных отношений, основанную на единстве высокоразвитых производительных сил и специфических капиталистических производственных отношений, в рамках которого машинное, фабрично-заводское и все более общественное по своему характеру производство соединяется с частной собственностью, наемным трудом, всеобщей товарностью, рыночной конкуренцией и анархией производства, вследствие чего само развитие капитализма осуществляется как расширенное воспроизводство капитала, сопровождающееся его накоплением, концентрацией и централизацией, воспроизводством классовой структуры общества, постоянным революционизированием техники, организации труда и производства, а также регулярным возникновением внутренних противоречий между общественным характером труда и частнокапиталистической формой присвоения, между безграничным стремлением к расширению и ограниченностью платежеспособного спроса, между необходимостью пропорционального воспроизводства и стихийной рыночной координацией.

Таким образом, капиталистическая экономика выступает специфической сферой промышленного цикла, поскольку именно в рамках нее имманентные противоречия капиталистического способа производства регулярно воспроизводят материальные основания для периодического развертывания, обострения и временного разрешения кризисов, а циклическая форма движения общественного воспроизводства получает свое наиболее полное и закономерное выражение.

\subsection{Капиталистическая экономика как внутренне противоречивая система}
Поскольку капиталистическая экономика представляет собой исторически конкретную форму организации общественного воспроизводства, ее анализ требует выделения содержащихся в самой ее основе фундаментальных противоречий, раскрывающих сущность, динамику и исторические пределы ее развития, задающих тем самым внутреннюю логику ее функционирования и обуславливающих закономерно возникающие в ходе ее движения диспропорции, кризисы и качественные переломы, формирующие основу предпосылок ее дальнейшего отрицания.

Материальной основой капиталистической экономики выступает исторически развернутый \textit{капиталистический способ производства}, содержащий в качестве основополагающих элементов \textit{товар}, \textit{деньги} и \textit{капитал}, каждый из которых заключает в себе собственные внутренние противоречия, однако наиболее глубокую форму принимают противоречия самого капиталистического способа производства как такового, обуславливающие нелинейный характер развития всей капиталистической экономики.

Таким образом, внутренне противоречивый характер капиталистической экономики обуславливает систему взаимосвязанных и иерархически организованных противоречий, среди которых \textit{главные противоречия} определяют общий способ движения, историческую направленность и пределы развития капиталистического строя, в то время как \textit{второстепенные противоречия} конкретизируют формы их проявления в отдельных процессах и сферах хозяйственной жизни, при этом характеризуясь исключительно временным разрешением в ходе регулярно возникающих диспропорций, кризисов и структурных переломов, которые не устраняют их окончательно, а всего лишь подготавливают почву для их дальнейшего воспроизводства на более высокой стадии развития капиталистической системы, подтверждая тем самым ключевую роль имманентных противоречий в качестве источника ее движения.

\subsection{Промышленный цикл в контексте фундаментальных закономерностей развития капиталистической экономики}
\textit{Промышленный цикл} представляет собой не внешнюю, эмпирически наблюдаемую периодичность колебаний хозяйственной активности, а специфическую форму проявления внутренней логики движения капиталистической экономики, отражающую ее имманентные противоречия и фундаментальные закономерности развития, а следовательно, требующую рассмотрения капиталистической системы не как механической совокупности отдельных явлений, а в качестве внутренне связанного, противоречивого и закономерно развивающегося целого, воспроизводящего собственные условия движения посредством последовательной смены фаз кризиса, депрессии, оживления и подъема.

Взаимно дополняющие друг друга процессы \textit{обобществления производства}, \textit{концентрации} и \textit{централизации} капитала, сопровождающиеся накоплением, технической модернизацией, ростом органического строения капитала, тенденцией к понижению нормы прибыли и относительным перенаселением, монополизацией экономики, вытеснением мелких и средних производителей, усилением рыночной власти крупнейших корпораций и деградацией конкурентной среды, формируют материальную основу периодических кризисов развития капитализма, выступающих в соответствии с законами диалектического и исторического материализма не случайными сбоями, а закономерными этапами движения капиталистической экономики, в рамках которых противоречия ее составных элементов и капиталистического способа производства в целом достигают предельной остроты, частично снимаются через разрушение избыточного капитала, перестройку воспроизводственных пропорций и усиление эксплуатации труда, после чего регулярно воспроизводятся вновь.

Таким образом, движение капиталистической экономики следует характеризовать именно в качестве \textit{промышленного цикла}, а не общей формы диалектического развития, поскольку опосредованные разрешением накопленных второстепенных противоречий качественные изменения в ходе присущих ей кризисов перепроизводства остаются сугубо в рамках капиталистического способа производства и не влекут за собой отрицание самой капиталистической системы как таковой, в то время как служащие первоисточником ее движения главные противоречия остаются незименными на всех стадиях ее развития. $\dagger$ (Приложение \ref{app:dialectica}, стр. \pageref{app:dialectica})